\documentclass[11pt,a4paper]{article}
\usepackage[margin=2.2cm]{geometry}
\usepackage[T1]{fontenc}
\usepackage[utf8]{inputenc}
\usepackage{lmodern}
\usepackage{setspace}
\usepackage{xcolor}
\usepackage{titlesec}

\definecolor{heading}{HTML}{2F5D50}
\titleformat{\section}{\Large\bfseries\color{heading}}{}{0pt}{}
\setlength{\parindent}{0pt}
\setlength{\parskip}{0.75em}
\onehalfspacing

\begin{document}

\begin{center}
    {\Huge\bfseries Freedom and Its Cost in \textit{West}}\\[0.5em]
    {\Large Analytical Essay}
\end{center}

\vspace{1em}

In the first quarter of Carys Davies' \textit{West}, freedom is shown as something attractive but also dangerous and unequal. John Cyrus Bellman leaves Pennsylvania because he believes that huge ancient creatures may still exist somewhere in the unexplored West. His journey seems free because he chooses it himself: he leaves his home, his sister Julie, and his daughter Bess in order to follow an idea that no one else believes in. However, the novel also shows that freedom is not simple. Bellman's choice gives him movement and purpose, but it also creates loss for the people he leaves behind. At the same time, Old Woman From A Distance's people are forced west, which creates a strong contrast between chosen movement and forced displacement. Through Bellman, Julie, Bess, and Old Woman From A Distance, \textit{West} suggests that freedom always depends on power, responsibility, and cost.

Bellman leaves behind the safety and duties of home. Most importantly, he leaves behind Bess, who must live without her father and wait for letters that may never come. He also leaves Julie, who must take responsibility for Bess while disagreeing with his decision. What he takes with him is much more than practical equipment. He takes supplies, trade goods, weapons, maps, and knowledge from the journals of Lewis and Clark, but he also takes an obsession. The newspaper report about the giant bones has changed the way he sees the world. For Bellman, the West becomes a place where imagination might become reality. This shows that his freedom is both physical and mental: he wants to move across the continent, but he also wants to prove that his dream is possible.

Bellman feels that he has to go because staying would mean living against his own imagination. Julie, Elmer, and others see the journey as foolish because it is dangerous, unrealistic, and irresponsible. Their criticism is reasonable, especially because Bellman is a father. Yet Bellman cannot simply ignore the idea of the enormous creatures. He is driven by curiosity and by the desire to see for himself. His freedom is therefore connected to individual desire. He believes that a meaningful life requires action, risk, and discovery. The problem is that his personal freedom is built on other people's sacrifice, especially Bess's loneliness.

Julie represents a very different kind of freedom. Unlike Bellman, she does not ride away into the unknown. Her life is connected to the home, to responsibility, and to practical limits. Bellman may see this as ordinary or restrictive, but the novel does not present Julie as weak. In many ways, she is more realistic than Bellman. Her freedom is limited by duty, yet she also has moral clarity. She understands that Bellman's choice affects more than himself. This contrast shows that freedom is not only the ability to leave; it can also mean choosing to stay, protect, and care for others.

The strongest contrast comes through Old Woman From A Distance. Bellman chooses to go west, but Old Woman's people are pushed west by pressure, unfair agreements, and the expansion of white settlers. For Bellman, the West represents possibility. For Old Woman's people, it represents loss: loss of land, safety, control, and future. This difference reveals that freedom in the novel is unequal. Some people have the power to treat the West as an adventure, while others experience it as exile. By placing Bellman's chosen journey beside Native displacement, Davies complicates the idea of frontier freedom. The same landscape can mean independence for one person and oppression for another.

This idea also connects to modern society. Today, freedom often means the ability to choose one's education, job, home, beliefs, or lifestyle. However, these choices are still shaped by money, family responsibilities, race, class, nationality, and political power. Some people can travel, take risks, or follow dreams because they have support and security. Others cannot, because their choices are limited by poverty, discrimination, or responsibility for others. Like Bellman, modern individuals may chase personal freedom while someone else pays the emotional or practical cost.

In conclusion, \textit{West} presents freedom as powerful but morally complicated. Bellman's journey shows the excitement of choosing one's own path, but it also shows the selfishness that can hide inside that choice. Julie and Bess reveal the cost of his freedom at home, while Old Woman From A Distance reveals the historical injustice behind westward movement. The novel therefore argues that freedom is never just about movement or personal desire. True freedom must include responsibility, fairness, and an awareness of who is forced to pay for another person's dream.

\end{document}
